\section{Tracciamento tridimensionale basato su modelli}

L'architettura di un sistema generico di stima e tracciamento basato su modello opera ricercando, all'interno dello spazio delle configurazioni, i parametri che minimizzano l'errore di corrispondenza tra le caratteristiche estratte dall'immagine in ingresso e quelle generate sinteticamente dal modello.
La struttura di tali sistemi si articola in quattro moduli: inizializzazione, predizione, estrazione delle caratteristiche e ricerca.

La fase di inizializzazione risulta strettamente necessaria al primo fotogramma, poiché, in assenza di stime precedenti, occorre definire la configurazione spaziale di partenza per avviare il tracciamento.
Tale problematica è frequentemente risolta ricorrendo a procedure manuali oppure chiedendo all'utente di assumere una posa predefinita e priva di occlusioni, come la mano completamente aperta.

Nei fotogrammi successivi interviene il modulo di predizione, il quale avvia la fase di ricerca a partire dalla posa precedente e basandosi su un modello di dinamica del sistema.

Il modulo dedicato all'estrazione delle caratteristiche ha il compito di isolare i dati visivi necessari al calcolo dell'errore di allineamento spaziale.
Nella maggior parte delle architetture vengono sfruttate caratteristiche bidimensionali sul piano immagine, richiedendo pertanto la proiezione continua del modello tridimensionale. Operazione computazionalmente onerosa.
Approcci alternativi e meno comuni impiegano caratteristiche tridimensionali acquisite tramite sensori di profondità o sistemi multicamera.

Infine, il modulo di ricerca esegue l'ottimizzazione vera e propria dello stato, valutando la somiglianza tra le metriche estratte e quelle sintetizzate dal modello.
Questo processo è storicamente implementato secondo due paradigmi principali. 
Il tracciamento a singola ipotesi si basa su una ricerca locale e mantiene unicamente la stima considerata migliore a ogni singola iterazione.
Esso risulta computazionalmente efficiente ma estremamente vulnerabile a perdite di tracciamento su sequenze prolungate.
I sistemi più avanzati adottano il tracciamento a ipotesi multiple per ovviare a queste criticità, il quale, mantenendo simultaneamente molteplici stime di posa, garantisce un'elevata tolleranza alle ambiguità visive e ai fallimenti locali, permettendo al processo di tracciamento di proseguire anche qualora l'ipotesi primaria dovesse rivelarsi inesatta.